\chapter{Introdução}
\label{cap:introducao}

A disciplina de \textbf{Fundamentos de Estruturas de Computadores}, ministrada pelo
Prof.\ Rubem Koide na Universidade Positivo, percorre um arco que vai dos
conjuntos numéricos e portas lógicas (aulas 2--4) até processadores
\textsc{cisc}/\textsc{risc} e sistemas operacionais (aulas 10--12),
passando pela arquitetura interna do computador e seus componentes
físicos (aulas 6--9). Em paralelo a essa progressão teórica, o curso
demanda dos alunos um projeto que materialize, em forma de artefato
didático, um recorte significativo desse conteúdo.

Este trabalho apresenta o desenvolvimento de um
\textbf{infográfico interativo sobre arquitetura de servidores},
escrito em \textsf{HTML5}, \textsf{CSS3} e \textsf{JavaScript ES6+}
sem qualquer dependência externa, executável diretamente em navegador
moderno. O recorte ``servidores'' foi escolhido por concentrar, num
único objeto de estudo, todos os subsistemas que a disciplina aborda:
do barramento e da hierarquia de memória até o escalonador de processos
e a virtualização --- portanto, exibe a teoria das aulas em ação.

\section{Justificativa}

A primeira versão do projeto, entregue em abril de 2026, recebeu como
crítica do orientador a \textit{baixa profundidade de conteúdo}: o
artefato apresentava analogias e curiosidades, mas pouca substância
técnica que se conectasse de fato com o material das aulas. Este
relatório descreve a \textbf{segunda iteração}, na qual o infográfico
foi reestruturado para incorporar:

\begin{itemize}
    \item tabelas comparativas com números reais de largura de banda,
          latência e custo (PCIe, UPI, DDR5, NVMe);
    \item fórmulas explícitas (Lei de Amdahl, cálculo de paridade
          em RAID 5, síndrome de Hamming \textsc{secded});
    \item \textit{specs} concretas de processadores citados em
          aula (Intel Core i7-5960X, Xeon Scalable, EPYC);
    \item simuladores executáveis (acesso à cache hit/miss,
          calculadora de IOPS, escalonador rubro-negro CFS); e
    \item bibliografia formal \textsc{abnt} \textsc{nbr 6023}
          com 17 referências, acessíveis diretamente da página
          via citações clicáveis.
\end{itemize}

\section{Objetivos}

\subsection{Objetivo geral}

Desenvolver um \textbf{infográfico interativo} que explique, de modo
visual e exploratório, os principais subsistemas que diferenciam um
servidor de um computador pessoal, conectando explicitamente cada
conceito ao conteúdo teórico das aulas 6--12 da disciplina.

\subsection{Objetivos específicos}

\begin{enumerate}
    \item Comparar Desktop \textit{vs.}\ Servidor em quatro dimensões
          --- alimentação, interconexão CPU--CPU, barramentos de E/S
          e gerenciamento remoto --- usando dados quantitativos.
    \item Demonstrar interativamente o comportamento de quatro
          níveis \textsc{raid} (0, 1, 5, 10) sob falha de disco e
          permitir o cálculo de capacidade útil e \textsc{iops} sob
          variação de parâmetros.
    \item Visualizar o código de Hamming \textsc{secded} usado
          em memória \textsc{ecc}, expondo a posição dos bits de
          paridade e o cálculo do síndrome.
    \item Apresentar a hierarquia de memória com latências em escala
          humana e simular a probabilidade de \textit{hit} em cada
          nível de cache.
    \item Explicar virtualização de Tipo~1 \textit{vs.}~Tipo~2,
          anéis de proteção, contêineres \textit{vs.}\ máquinas
          virtuais e o escalonador CFS do Linux.
    \item Integrar todas as camadas anteriores numa visão única do
          \textit{stack} completo de um servidor moderno.
\end{enumerate}

\section{Estrutura do trabalho}

O Capítulo~\ref{cap:referencial} estabelece o referencial teórico,
agrupado pelos seis subsistemas tratados no infográfico. O
Capítulo~\ref{cap:metodologia} descreve materiais, métodos e o
ambiente de desenvolvimento, incluindo as ferramentas de IA assistiva
empregadas. O Capítulo~\ref{cap:desenvolvimento} apresenta o
infográfico módulo a módulo, com tabelas, figuras e descrição das
simulações implementadas. O Capítulo~\ref{cap:conclusao} sintetiza os
resultados, limitações e trabalhos futuros. O Anexo~A traz o
\emph{link} do código-fonte completo no GitHub e trechos críticos do
HTML/JavaScript desenvolvido.
